\begin{frame}{유명한 사례: 두 번 발견된 알고리즘}
  k-means의 기원 = ML 역사에서 \textbf{``동시 독립
  발견''}(simultaneous discovery)의 대표 사례:
  \vskip0.6em
  \begin{columns}[T]
    \begin{column}{0.48\textwidth}
      \kb{로이드 (Lloyd, 1957)}
      \begin{itemize}\small
        \item 벨 연구소, \textbf{PCM}(상용 전화 신호를 디지털로
          변환)의 양자화 문제 연구 중 발명
        \item 정식 논문 전에 \textbf{내부 보고서}로만 머물러 있었음
      \end{itemize}
    \end{column}
    \begin{column}{0.48\textwidth}
      \kb{매퀸 (MacQueen, 1967)}
      \begin{itemize}\small
        \item 스탠퍼드, \textbf{클러스터링} 문제에서
          \textbf{독립적으로} 같은 알고리즘 재발명
        \item ``k-means''라는 \textbf{이름}과 함께 발표
        \item 로이드가 훨씬 먼저였다는 사실이 알려지기까지 꽤
          오래 걸림
      \end{itemize}
    \end{column}
  \end{columns}
  \vskip0.6em
  두 발견은 \textbf{문제가 달랐는데}(신호 양자화 vs 데이터
  클러스터링), 핵심 구조(\textbf{거리로 가장 가까운 대표에 할당
  $\to$ 대표 재계산})가 정확히 같았다.
  \vskip0.4em
  \begin{block}{교훈}
    \textbf{``거리 기반 반복 재할당''은 특정 도메인의 트릭이
    아니라, 거리 함수 위에서 자연스러운 최적화 구조 그
    자체} --- 같은 패턴이 서로 다른 문제에서 반복 등장하는 이유.
  \end{block}
\end{frame}
